\section{Evaluation: What Holds vs.\ What Breaks}
\label{sec:eval}

% The headline payoff. A LIKE-FOR-LIKE comparison on ONE substrate under ONE
% threat model (Sec. 2): the same attack suite, same metrics, same model — run
% against static obfuscation (AloePri, Sec. 5) and the TEE+GPU split (Sec. 4).
% This table IS the paper's thesis in one place.

\subsection{Comparison Basis and Setup}
% LIKE-FOR-LIKE = same adversary (HbC, white-box weights+scheme), same goal
% (recover input tokens), same attack suite, same base model. The OBSERVATION
% SURFACE varies BY DESIGN: AloePri exposes obfuscated I/O+weights; the TEE-split
% exposes obfuscated activations. The surface IS the defense, so it is the
% independent variable -- not a confound. This is NOT "same observed bytes"; do
% not claim "single substrate" (corrects stage-1 phrasing). Attacker budget is
% pinned to each scheme's own claimed regime (\cref{sec:threat}).
\needswork{state the comparison basis explicitly; then model, hardware, attack
harness, metrics (TTRSR/CosSim/BLEU), datasets}

\subsection{Head-to-Head: Attack Resistance}
\begin{table*}[t]
  \centering
  \caption{Attack success by defense, under the same adversary, target, and
           attack suite (observation surface varies by architecture).
           Lower is safer. The headline contrast.}
  \label{tab:holds-vs-breaks}
  \begin{tabular}{@{}lcccc@{}}
    \toprule
    Defense & Inversion & GRAM & ISA AttnScore & Overhead \\
    \midrule
    None (plaintext)      & \needswork{} & \needswork{} & \needswork{} & --- \\
    Static obf.\ (AloePri)& \needswork{high} & \needswork{high} & \needswork{high} & low \\
    TEE+GPU split (ours)  & \needswork{low}  & \needswork{low}  & \needswork{low}  & \needswork{} \\
    \bottomrule
  \end{tabular}
\end{table*}

\subsection{The Defense/Accuracy Frontier}
% Where static obfuscation sits on the defense-vs-accuracy curve, and why the
% split escapes that frontier (at an overhead cost).
\needswork{frontier plot + discussion}

\subsection{Takeaways}
\needswork{state the holds/breaks conclusion plainly; bound it by the threat model}
